\documentclass[12pt]{article}
%^ Start of document
\usepackage{times}  %Makes text TimesNewRoman
\usepackage{setspace} %Allows you to set single or double space 
\usepackage[
sorting=none
]{biblatex}
%\usepackage{biblatex} %Imports biblatex package
\usepackage{graphicx} % Required for inserting images
\usepackage{authblk}
\usepackage[colorlinks=true, urlcolor=blue, linkcolor=blue]{hyperref}
\usepackage{subfig}
\usepackage{float}
\usepackage{tabularx}
\usepackage{multirow}
\usepackage[paperheight=11in,
   paperwidth=8.5in,
   top=20mm,
   bottom=20mm,
   left=40mm,
   right=40mm]{geometry}
\usepackage{moresize}
\usepackage{tikz}
\usepackage{xcolor}
\usepackage{physics}
\usepackage{amssymb}
\usepackage{mathrsfs}
\usepackage{amsmath}
\usepackage[document]{ragged2e}
\usepackage{comment}
\usepackage{xcolor}
\newcommand*{\SignatureAndDate}[3]{
    \par\noindent\makebox[2.5in]{\hrulefill} \hfill\makebox[2.0in]{\hrulefill}
    \par\noindent\makebox[2.5in][l]{#1}      \hfill\makebox[2.0in][l]{Date}
    \par\noindent\makebox[2.5in][l]{#2}
    \par\noindent\makebox[2.5in][l]{#3}}

    
\addbibresource{mm.bib}
\begin{document}
\singlespacing  %Sets single spacing for whole document
\title{Suppression of radiation emissions from relativistic plasma by utilizing magnetic perturbations}

\author[1]{Zachary R. Tullier}
\affil[1]{Student\\Department of Physical and Applied Science\\ University of Houston - Clear Lake \\Houston, Texas\\U.S}

\date{}
\maketitle

\begin{center}
    {\fontfamily{qcr}\selectfont APPROVED BY THE THESIS COMMITTEE}
\end{center}

\SignatureAndDate{Samina S. Masood}{Ph.D., Chair}{Thesis Advisor}
\vspace{.2in}
\SignatureAndDate{Serkan Caliskan}{Ph.D}{Thesis Committee}
\vspace{.2in}
\SignatureAndDate{Brandon Reddell}{Ph.D}{Thesis Committee}

\vspace{.3in}
\begin{center}
    {\fontfamily{qcr}\selectfont APPROVED BY THE COLLEGE OF SCIENCE AND ENGINEERING}
\end{center}

\SignatureAndDate{David Garrison}{Ph.D}{Associate Dean}
\vspace{.2in}
\SignatureAndDate{Jennifer Irvin}{Ph.D}{Dean}

\newpage

\begin{abstract}
    In atomic systems, magnetic moments originate from both the orbital motion and intrinsic spin of electrons. These magnetic properties play a significant role in plasma behavior, especially in highly magnetized environments such as those found in fusion reactors. One of the major challenges in achieving efficient fusion is minimizing energy losses, particularly those due to bremsstrahlung radiation, which is often intensified by superthermal electron populations. This thesis proposes to investigate the application of controlled magnetic perturbations in low-energy fusion plasmas as a method for suppressing bremsstrahlung losses. By exploring the interaction between externally imposed magnetic field structures and electron dynamics, the study aims to identify regimes where such perturbations reduce high-energy X-ray emissions and improve overall energy confinement. The findings are expected to inform more effective magnetic control strategies for next-generation fusion devices.
\end{abstract}

\section{Introduction: Bremsstrahlung losses in Relativistic $p^{-11}B$ Fusion Plasmas}

Nuclear fusion is a promising pathway toward clean and sustainable energy. Among the various fusion reactions, proton–boron-11 (p–¹¹B) fusion stands out as an attractive aneutronic option, producing minimal neutron radiation and potentially enabling direct energy conversion. In addition, its fuel sources are abundant and could even be synthesized from nuclear fission waste, adding to its long-term appeal.
However, the practical realization of p–¹¹B fusion faces significant challenges, the most critical of which is radiative energy loss. In high-temperature plasmas, energetic electrons emit bremsstrahlung X-rays when scattered by ionic Coulomb fields. This radiation loss can exceed the energy generated by fusion reactions, especially in plasmas with relativistic electron populations (\cite{SuppressionBremsstrahlung}). At these energies, bremsstrahlung emission is further intensified by the high-energy tail of the electron distribution, and electron–electron collisions also contribute significantly—unlike in lower-temperature plasmas.
This fundamental energy loss mechanism severely limits the net energy gain from p–¹¹B fusion and threatens its feasibility as a practical energy source. Therefore, identifying methods to reduce or suppress bremsstrahlung radiation is a critical step toward enabling viable aneutronic fusion. This research aims to explore magnetic perturbation as a potential technique to mitigate such losses, thereby contributing to the broader effort to make p–¹¹B fusion a viable energy solution.

\subsection{Relativistic Bremsstrahlung Theory in Strong Fields}

At relativistic energies, the bremsstrahlung differential cross-section (Bethe–Heitler formula, see \cite{Bethe-Heitler}) grows with electron energy, meaning superthermal electrons dominate X-ray emission (\cite{SuppressionBremsstrahlung}). This has two important implications for suppression strategies. First, shaping the electron energy distribution to minimize the high-energy tail can disproportionately cut radiative losses (\cite{SuppressionBremsstrahlung}). Second, intense electromagnetic fields can alter electron trajectories and quantum states, modifying emission probabilities. For example, in ultra-strong magnetic fields, electrons occupy discrete Landau energy levels (quantized transverse motion), which changes collision dynamics. If the field is extremely high (above $~5×10^9$G), ion–electron collisions can no longer easily excite electrons to higher Landau levels (\cite{SuppressionBremsstrahlung}, \cite{ProspectsP11B}). In effect, Landau quantization raises the threshold for energy transfer to electrons, keeping electron temperatures lower relative to ions. This can dramatically reduce bremsstrahlung cooling in principle, as theorized for magnetized proton–boron plasmas where fusion power could far exceed X-ray losses under multi-gigagauss fields (\cite{ProspectsP11B}). Such quantum magnetic-field effects (Landau quantization and related QED phenomena) form an extreme end of bremsstrahlung suppression – likely relevant in astrophysical or dense plasma focus scenarios, but requiring magnetic fields beyond conventional lab scales.

Strong magnetic fields also introduce classical suppression effects. In relativistically intense laser–plasma interactions, quasi-static magnetic fields on the order of tens of megatesla are predicted to reduce bremsstrahlung emission probability by altering electron orbits (\cite{10.1063/5.0167288}). This magnetic suppression mechanism occurs because a large magnetic field bends electron trajectories, effectively limiting the duration and straight-line distance over which an accelerating electron can emit a bremsstrahlung photon. A study implemented this effect in Particle-in-Cell (PIC) simulations using the EPOCH code, incorporating magnetic and electric field influences into the bremsstrahlung emission module. Their study, the first of its kind, showed that a strong field suppresses especially the low-energy X-ray emissions and also alters the dynamics of radiating electrons. This suggests that in magnetized relativistic plasmas, bremsstrahlung can be partially mitigated by field-induced trajectory perturbations, complementing any quantum Landau-level effects.

\subsection{Tailoring Electron Distributions (Energy Cutoff Method)}

Because high-energy electrons disproportionately drive bremsstrahlung, one intuitive approach is to remove or redistribute the electron tail to lower energies. This idea can be realized by introducing an energy cutoff in the electron distribution. Instead of a pure Maxwell–Jüttner distribution, the plasma electrons follow a distribution where beyond some cutoff energy $E_c$ the population is sharply reduced (\cite{SuppressionBremsstrahlung}). Recent simulations and analytic models have explored how such cutoffs suppress radiation. The study by Munirov and Fisch (\cite{Munirov_2023}) showed that redistributing superthermal electrons into lower energies can significantly reduce bremsstrahlung losses in a relativistic plasma. Importantly, this suppression effect is a uniquely relativistic benefit, in nonrelativistic plasmas removing the electron tail has much less impact on total radiation. Munirov and Fisch calculated bremsstrahlung power for a range of cutoff energies and electron temperatures, keeping total electron density fixed (i.e. electrons are not removed from the system but slowed down). They found that even a modest cutoff can yield measurable reductions. For example, at an electron temperature $T_e = 100$keV (typical of p–¹¹B reactors), imposing a cutoff around $E_c \approx 70-100$keV can lower electron–ion bremsstrahlung power by ~10\% and electron–electron bremsstrahlung by ~20\% relative to a Maxwellian plasma. Higher cutoffs (>2–3 times the thermal energy) see diminishing returns, as most electrons remain below the cutoff.

There are several physical mechanisms to achieve such an effective energy cutoff. In open magnetic field-line devices (like mirror machines or certain inertial-electrostatic confinement setups), high-energy electrons can be allowed to escape along field lines, while colder electrons and ions are confined by electrostatic potentials. The escaping hot electrons carry kinetic energy out, effectively truncating the distribution inside the core plasma. As one example, a magnetic mirror or ambipolar electrostatic well can be tuned so that electrons above a certain energy have enough parallel velocity to overcome the confining potential and stream out, whereas lower-energy electrons remain trapped. This selective loss mechanism continuously removes the nascent fast electrons as they are produced by collisions, preventing a high-energy tail from accumulating. The lost electrons’ charge is compensated by ion outflow or by replenishing with slower electrons to maintain quasineutrality. Overall, the plasma sustains a skewed distribution with far fewer relativistic electrons than a Maxwellian, thereby emitting less bremsstrahlung (\cite{SuppressionBremsstrahlung}).

Another approach to distribution control is through waves and instabilities that damp high-energy electrons. Resonant cyclotron emission can act as a radiative sink for electrons above a threshold energy, as noted in strong-field plasma calculations. In a 1999 study (\cite{emissivity}), Riffert found that in a magnetized proton–electron plasma, electrons above the cyclotron resonance energy rapidly lose energy to intense cyclotron (synchrotron) radiation, depleting the high-energy tail and leaving a narrow spectral line at the cyclotron frequency. This indicates that under certain conditions, synchrotron cooling can outpace bremsstrahlung for the fastest electrons, effectively capping their energy and reducing continuum X-ray emission. Similarly, modern PIC simulations of laser-driven plasmas have observed that including radiation reaction forces (synchrotron losses) leads to a natural high-energy cutoff in the electron spectrum. While synchrotron emission is itself an energy loss channel, it tends to be more peaked (narrow-band) and may be easier to capture or mitigate than broadband bremsstrahlung. This interplay of radiation processes offers multiple avenues to engineer the electron distribution for minimal bremsstrahlung: either by external fields that allow fast electron escape or by internal plasma processes that selectively drain energy from fast electrons.

\subsection{Resonant Magnetic Perturbations and Stochastic Confinement}

Magnetic perturbation techniques exploit three-dimensional fields to alter confinement, and they can be a powerful tool for controlling high-energy particles. In toroidal fusion devices (tokamaks and stellarators), applying resonant magnetic perturbations (RMPs) has been well-studied as a method to induce magnetic islands and stochastic regions in the plasma edge. These perturbations, typically small static fields introduced by external coils, resonate with the plasma’s field geometry and partially break the confining magnetic surfaces. The resulting magnetic islands and open field lines can allow particles (and heat) to flow out of the plasma more readily. While RMPs are traditionally used for edge-localized mode (ELM) control, recent work indicates they can also enhance the loss of runaway electrons and other super-thermal particles, directly relevant to bremsstrahlung suppression since runaways generate intense X-ray emission.

Simulation studies have examined how different RMP spectra affect fast electron confinement. Cornille (\cite{10.1063/5.0087314}) performed MHD and test-particle simulations of runaway electron behavior in the MST tokamak with applied RMPs. Using the NIMROD code (for nonlinear MHD) coupled with tracer electrons (using the KORC orbit-following code), they found that an RMP with a certain mode structure can chaoticize the magnetic field in the plasma edge and substantially deconfine high-energy electrons. In their results, an $m=3$ RMP (with three island chains around the torus) induced a stochastic outer region that led to ~50\% of runaway electrons being lost from the plasma, compared to only ~10–15\% loss without RMP or with a non-resonant $m=1$ perturbation. The RMP effectively opened up escape paths for the energetic electrons, preventing a fully developed runaway beam. These findings are consistent with experimental observations on devices like J-TEXT and TEXTOR, where applying resonant $n=1$ or $n=2$ perturbation fields during disruptions measurably increased runaway electron losses (sometimes achieving full suppression of a runaway current). In the fusion plasma context, RMPs introduce a controllable level of magnetic chaos: by tuning the perturbation amplitude and spectrum, one can strike a balance between purging high-energy electrons and still confining the bulk plasma. This concept directly parallels the energy-cutoff approach, instead of relying on ambipolar electrostatic fields as in mirror machines, a tokamak could use magnetic topological tweaks to let the most energetic electrons escape (or hit a limiter) before they radiate.

State-of-the-art simulation tools support these studies of 3D field effects. The JOREK non-linear MHD code, for instance, has been used to model plasma response to RMPs in tokamaks, capturing the formation of magnetic islands and the resultant changes in transport. Similar codes like NIMROD and M3D-C1 can include drift-kinetic or fluid models for runaway electrons, allowing integrated simulations of how applied perturbations dissipate fast electron populations. Meanwhile, HINT/HINT2 (an MHD equilibrium solver) and analytic field-line tracing codes like FLARE provide ways to visualize and design magnetic configurations with open field lines (\cite{FLARE}). FLARE in particular can compute Poincaré maps and connection lengths in non-axisymmetric fields, helping researchers identify how an RMP or helical field will intersect material walls or direct particles to a divertor target. Using these tools, one can simulate a scenario where, for example, slow electrons remain confined on closed flux surfaces while fast electrons wander into stochastic regions and are removed, effectively creating the desired energy-selective confinement. Such magnetic perturbation-based techniques are a promising avenue for bremsstrahlung suppression in advanced p–¹¹B reactors, especially in designs that start from a toroidal confinement concept and need to shed excess electron energy.

\section{Thesis Objective}
    This proposed thesis will provide an in-depth review of novel bremsstrahlung suppression techniques that utilize state-of-the-art simulation tools to model resonant magnetic perturbations. A simple modeled using these same techniques will be created to evaluate the lower limit of feasibility using instrumentation that is currently in service.

\section{Potential Implications}
    There has been a dramatic increase in funding to build, test, and operate magnetic confinement fusion energy plants over the past decade. These efforts, driven by public and private investment alike, underscore a global push toward realizing practical fusion power. Within this context, aneutronic fusion systems—such as proton–boron-11 (p–¹¹B) reactors—offer significant long-term advantages, including minimal radioactive waste, reduced shielding requirements, and the possibility of direct energy conversion. The results of this thesis could provide a pathway for bremsstrahlung suppression techniques that improve the feasibility of p–¹¹B fusion and related advanced fuel cycles.

    One key implication of this research is the potential to reduce the engineering complexity and cost of future fusion reactors. Bremsstrahlung losses impose strict thermal and radiative constraints on the design of reactor components, particularly the first wall and divertor. By minimizing X-ray output from relativistic electrons, magnetic perturbation-based suppression methods could reduce the required thickness of radiation shielding materials and cooling systems. This, in turn, can accelerate the deployment timeline for commercial fusion systems and lower capital costs, making fusion more economically competitive with other energy sources.

    In addition, successful suppression of bremsstrahlung emissions may enable novel confinement regimes and plasma configurations previously considered unfeasible. For example, current designs for aneutronic fusion rely on extreme magnetic field strengths or complex non-Maxwellian electron distributions to remain net-power positive. If magnetic perturbations can selectively remove or deconfine superthermal electrons, it could open the door for new operational scenarios at lower magnetic field strengths and more manageable plasma conditions. This could benefit a wide range of confinement approaches, including spherical tokamaks, field-reversed configurations, and compact toroidal reactors.

    Furthermore, the modeling techniques developed and refined in this study—such as MHD simulations incorporating resonant magnetic perturbations, particle orbit tracing, and bremsstrahlung diagnostics—will contribute to the broader computational plasma physics community. These tools can be extended to study runaway electron mitigation in deuterium–tritium (D–T) reactors, laser–plasma interactions, and astrophysical systems, where radiative energy loss also plays a dominant role. The results could also help inform safety analyses for pulsed operation or disruption scenarios in experimental devices.

    Finally, from a policy and planning perspective, identifying magnetic confinement concepts that can operate effectively with reduced bremsstrahlung losses may influence funding strategies and experimental priorities. If this research establishes clear theoretical or simulated thresholds for effective radiation suppression using existing magnetic field technologies, it could help determine which experimental fusion programs are most likely to benefit from aneutronic fuel development. This includes evaluating international initiatives like TAE Technologies, HB11 Energy, and Dense Plasma Focus devices, as well as government-backed facilities such as SPARC, ITER, or DEMO.

\printbibliography    
\end{document}
